% Flowchart transcription from img/2010/new_curriculum_liberal/q8_options.png.
% Nodes: 开始; 输入 N; k=1,S=0; S=S+1/[k(k+1)]; k<N?; k=k+1; 输出 S; 结束.
% Directed edges: 开始→输入→初始化→求和→判定; 判定(否)→输出→结束;
% 判定(是)→k递增→上方回线→求和.
\begin{tikzpicture}[
  >=Stealth,
  line width=.55pt,
  line cap=round,
  line join=round,
  every node/.style={font=\normalsize,anchor=center,align=center,
    execute at begin node={\setlength{\parindent}{0pt}\noindent}},
  flowbox/.style={draw,minimum height=.70cm,inner xsep=4pt,inner ysep=2pt,
    anchor=center,align=center},
  terminal/.style={flowbox,rounded corners=.18cm,minimum width=1.42cm,
    minimum height=.76cm},
  io/.style={flowbox,shape=trapezium,trapezium left angle=72,
    trapezium right angle=108,minimum height=.82cm},
  decision/.style={flowbox,shape=diamond,aspect=2.75,inner sep=2pt},
  flowarrow/.style={->,line width=.55pt},
]
  \node[terminal] (start) at (0,12.10) {开始};
  \node[io,minimum width=3.35cm] (input) at (0,10.62) {输入 $N$};
  \node[flowbox,minimum width=4.35cm] (init) at (0,9.05)
    {$k=1,\ S=0$};
  \node[flowbox,minimum width=4.55cm,minimum height=1.08cm] (sum) at (0,7.05)
    {$S=S+\dfrac{1}{k(k+1)}$};
  \node[decision,minimum width=4.25cm,minimum height=1.08cm] (test) at (0,5.20)
    {$k<N$};
  \node[flowbox,minimum width=2.85cm] (inc) at (4.35,7.30)
    {$k=k+1$};
  \node[io,minimum width=3.15cm] (output) at (0,3.15) {输出 $S$};
  \node[terminal] (finish) at (0,1.48) {结束};

  \coordinate (loopjoin) at (0,8.20);
  \coordinate (branchright) at (4.35,5.20);
  \coordinate (righttop) at (4.35,8.20);
  \draw[flowarrow] (start.south) -- (input.north);
  \draw[flowarrow] (input.south) -- (init.north);
  \draw (init.south) -- (loopjoin);
  \draw[flowarrow] (loopjoin) -- (sum.north);
  \draw[flowarrow] (sum.south) -- (test.north);
  \draw[flowarrow] (test.south) -- node[right=3pt] {否} (output.north);
  \draw[flowarrow] (output.south) -- (finish.north);
  \draw (test.east) -- node[above=2pt] {是} (branchright);
  \draw[flowarrow] (branchright) -- (inc.south);
  \draw[flowarrow] (inc.north) -- (righttop) -- (loopjoin);
  \path[use as bounding box] (-2.92,1.05) rectangle (5.70,12.48);
\end{tikzpicture}
